% genere a la main : annexe des tables de benchmark completes.
% Les fragments viennent de work/final_tables_and_plots/tables/, francises par
% make_annexe_tables.py qui en depose des copies dans figs/tables/.
% Chaque table par modele occupe une page ; d'ou leur renvoi en annexe.

\chapter{Tables de benchmark complètes}
\label{ann:tables}

Cette annexe rassemble les mesures détaillées de l'axe mono-TPU, trop
volumineuses pour figurer dans le corps du chapitre~\ref{chap:resultats}. Les
tables~\ref{tab:baseline-speedup} et~\ref{tab:quant-speedup} portent sur les
modèles de référence non élagués ; les suivantes donnent, pour chacune des 7
architectures, le détail par taux d'élagage et par critère d'importance.

Les colonnes reprennent les grandeurs définies en section~\ref{sec:metriques} :
précision INT8 mesurée sur l'accélérateur et son écart à la référence, latence
Edge TPU, accélérations rapportées au modèle non élagué et au processeur,
réductions réelles du nombre de paramètres et des opérations, et répartition du
modèle entre mémoire interne et mémoire externe telle que la rapporte le
compilateur.

\section{Modèles de référence}

{\footnotesize
\begin{table}[H]
  \centering
  \caption{Latence des modèles de référence sur CPU et sur Edge TPU en INT8. Les valeurs sont des moyennes $\pm$ 1$\sigma$ sur $N{=}200$ inférences après chauffe ; l'incertitude sur l'accélération est propagée.}
  \label{tab:baseline-speedup}
  \centerline{%
  \begin{tabular}{l r rrr}
    \toprule
    Modèle & Taille (Mio) & CPU INT8 (ms) & TPU INT8 (ms) & Accél. TPU/CPU \\
    \midrule
    ResNet-18      & 10,84 & $74{,}37 \pm 0{,}76$  & $11{,}45 \pm 0{,}02$ & $6{,}50\times \pm 0{,}07$ \\
    ResNet-50      & 23,30 & $190{,}16 \pm 5{,}97$ & $50{,}92 \pm 0{,}49$ & $3{,}73\times \pm 0{,}12$ \\
    VGG-19         & 19,30 & $69{,}87 \pm 3{,}48$  & $39{,}29 \pm 0{,}47$ & $1{,}78\times \pm 0{,}09$ \\
    WRN-28-10      & 35,13 & $677{,}81 \pm 8{,}98$ & $92{,}13 \pm 0{,}38$ & $7{,}36\times \pm 0{,}10$ \\
    MobileNetV2    &  2,70 & $9{,}67 \pm 0{,}09$   & $1{,}24 \pm 0{,}01$  & $7{,}79\times \pm 0{,}08$ \\
    GoogLeNet      &  5,66 & $63{,}02 \pm 0{,}52$  & $2{,}04 \pm 0{,}01$  & $30{,}90\times \pm 0{,}33$ \\
    SqueezeNet 1.1 &  0,83 & $4{,}61 \pm 0{,}15$   & $0{,}67 \pm 0{,}01$  & $6{,}91\times \pm 0{,}27$ \\
    \bottomrule
  \end{tabular}}
\end{table}

\begin{table}[H]
  \centering
  \caption{Effet de la quantification sur CPU pour les modèles de référence, du flottant vers l'INT8. Les latences sont des moyennes $\pm$ 1$\sigma$ sur $N{=}200$ inférences après chauffe ; l'incertitude sur l'accélération est propagée.}
  \label{tab:quant-speedup}
  \centerline{%
  \begin{tabular}{l rrrr}
    \toprule
    Modèle & CPU FP32 (ms) & CPU INT8 (ms) & Accél. INT8/FP32 & $\Delta$ Top-1 \\
    \midrule
    ResNet-18      & $117{,}66 \pm 3{,}12$  & $74{,}37 \pm 0{,}76$  & $1{,}58\times \pm 0{,}04$ & $-0{,}33$ \\
    ResNet-50      & $306{,}36 \pm 4{,}95$  & $190{,}16 \pm 5{,}97$ & $1{,}61\times \pm 0{,}06$ & $-0{,}38$ \\
    VGG-19         & $147{,}72 \pm 2{,}65$  & $69{,}87 \pm 3{,}48$  & $2{,}11\times \pm 0{,}11$ & $-0{,}36$ \\
    WRN-28-10      & $703{,}22 \pm 12{,}41$ & $677{,}81 \pm 8{,}98$ & $1{,}04\times \pm 0{,}02$ & $-0{,}66$ \\
    MobileNetV2    & $17{,}32 \pm 1{,}23$   & $9{,}67 \pm 0{,}09$   & $1{,}79\times \pm 0{,}13$ & $-2{,}48$ \\
    GoogLeNet      & $100{,}01 \pm 2{,}55$  & $63{,}02 \pm 0{,}52$  & $1{,}59\times \pm 0{,}04$ & $-0{,}46$ \\
    SqueezeNet 1.1 & $7{,}38 \pm 0{,}63$    & $4{,}61 \pm 0{,}15$   & $1{,}60\times \pm 0{,}15$ & $-5{,}29$ \\
    \bottomrule
  \end{tabular}}
\end{table}

}

\clearpage
\section{Détail par architecture}

\renewcommand{\benchscale}{0.70}
\begin{table}[H]
  \centering
  \caption{Mesures complètes de ResNet-18 élagué, en INT8.}
  \label{tab:bench-resnet18}
  \centerline{\resizebox{\benchscale\linewidth}{!}{%
  \begin{tabular}{c l rr rrr rr rrr}
    \toprule
    & & \multicolumn{2}{c}{Précision} & \multicolumn{3}{c}{Temps} & \multicolumn{2}{c}{Réduction (\%)} & \multicolumn{3}{c}{Mémoire (Mio)} \\
    \cmidrule(lr){3-4} \cmidrule(lr){5-7} \cmidrule(lr){8-9} \cmidrule(lr){10-12}
    Cible & Critère & Top-1 & $\Delta$ Top-1 & Lat. (ms) & Accél. TPU & Accél. CPU & Param. & MACs & Taille & Interne & Externe \\
    \midrule
    -- & référence & 76,55 & +0,00 & 11,45 & 1,00$\times$ & 1,00$\times$ & 0,0 & 0,0 & 10,84 & 7,65 & 3,11 \\
    \midrule
    \multirow{7}{*}{10\%} & Magnitude L1 & 76,67 & +0,12 & 9,50 & 1,20$\times$ & 1,49$\times$ & 7,7 & 37,7 & 10,00 & 7,65 & 2,51 \\
     & Magnitude L2 & 76,46 & -0,09 & 9,52 & 1,20$\times$ & 1,48$\times$ & 7,7 & 37,4 & 10,00 & 7,65 & 2,51 \\
     & BN-Scale & 76,83 & +0,28 & 9,58 & 1,19$\times$ & 1,38$\times$ & 10,0 & 32,5 & 9,75 & 7,65 & 2,53 \\
     & FPGM & 76,57 & +0,02 & 9,82 & 1,17$\times$ & 1,42$\times$ & 10,0 & 34,5 & 9,75 & 7,65 & 2,61 \\
     & Taylor & 77,04 & +0,49 & 9,62 & 1,19$\times$ & 1,34$\times$ & 10,0 & 30,0 & 9,75 & 7,65 & 2,55 \\
     & OBD-C & 76,84 & +0,29 & 9,61 & 1,19$\times$ & 1,34$\times$ & 10,0 & 29,5 & 9,75 & 7,65 & 2,54 \\
     & Random & 76,57 & +0,02 & 10,03 & 1,14$\times$ & 1,07$\times$ & 10,1 & 11,0 & 9,75 & 7,64 & 2,67 \\
    \midrule
    \multirow{7}{*}{20\%} & Magnitude L1 & 75,90 & -0,65 & 6,53 & 1,75$\times$ & 1,98$\times$ & 19,6 & 55,2 & 8,72 & 7,69 & 1,59 \\
     & Magnitude L2 & 75,64 & -0,91 & 6,52 & 1,76$\times$ & 2,04$\times$ & 19,9 & 55,4 & 8,69 & 7,65 & 1,60 \\
     & BN-Scale & 76,08 & -0,47 & 6,99 & 1,64$\times$ & 1,73$\times$ & 20,1 & 47,5 & 8,67 & 7,65 & 1,74 \\
     & FPGM & 76,48 & -0,07 & 5,98 & 1,91$\times$ & 1,70$\times$ & 20,2 & 46,0 & 8,66 & 7,65 & 1,44 \\
     & Taylor & 76,56 & +0,01 & 6,30 & 1,82$\times$ & 1,47$\times$ & 20,1 & 38,0 & 8,67 & 7,65 & 1,53 \\
     & OBD-C & 76,57 & +0,02 & 6,97 & 1,64$\times$ & 1,49$\times$ & 20,1 & 37,3 & 8,67 & 7,65 & 1,74 \\
     & Random & 76,49 & -0,06 & 7,96 & 1,44$\times$ & 1,20$\times$ & 20,1 & 21,1 & 8,67 & 7,65 & 2,02 \\
    \midrule
    \multirow{7}{*}{30\%} & Magnitude L1 & 74,46 & -2,09 & 3,42 & 3,34$\times$ & 2,53$\times$ & 30,0 & 66,3 & 7,59 & 7,73 & 0,70 \\
     & Magnitude L2 & 74,92 & -1,63 & 2,25 & 5,09$\times$ & 2,41$\times$ & 30,1 & 64,5 & 7,59 & 7,69 & 0,40 \\
     & BN-Scale & 75,87 & -0,68 & 3,60 & 3,18$\times$ & 1,97$\times$ & 30,1 & 55,2 & 7,59 & 7,65 & 0,76 \\
     & FPGM & 75,51 & -1,04 & 1,94 & 5,90$\times$ & 2,10$\times$ & 30,0 & 57,5 & 7,59 & 7,65 & 0,30 \\
     & Taylor & 76,29 & -0,26 & 2,81 & 4,08$\times$ & 1,61$\times$ & 30,1 & 43,8 & 7,60 & 7,65 & 0,54 \\
     & OBD-C & 76,55 & +0,00 & 3,57 & 3,20$\times$ & 1,69$\times$ & 30,1 & 44,5 & 7,60 & 7,65 & 0,73 \\
     & Random & 76,45 & -0,10 & 2,81 & 4,07$\times$ & 1,38$\times$ & 30,0 & 30,8 & 7,60 & 7,65 & 0,54 \\
    \midrule
    \multirow{7}{*}{40\%} & Magnitude L1 & 72,96 & -3,59 & 1,07 & 10,72$\times$ & 3,20$\times$ & 40,0 & 74,2 & 6,51 & 6,91 & 0,00 \\
     & Magnitude L2 & 73,66 & -2,89 & 1,10 & 10,37$\times$ & 2,93$\times$ & 40,1 & 71,5 & 6,51 & 7,15 & 0,00 \\
     & BN-Scale & 74,70 & -1,85 & 1,19 & 9,64$\times$ & 2,48$\times$ & 40,1 & 64,3 & 6,51 & 7,07 & 0,00 \\
     & FPGM & 75,25 & -1,30 & 1,09 & 10,53$\times$ & 2,72$\times$ & 40,1 & 68,6 & 6,51 & 6,79 & 0,00 \\
     & Taylor & 76,48 & -0,07 & 1,37 & 8,37$\times$ & 1,73$\times$ & 40,0 & 47,0 & 6,53 & 7,06 & 0,00 \\
     & OBD-C & 75,97 & -0,58 & 1,32 & 8,69$\times$ & 1,82$\times$ & 40,1 & 49,8 & 6,52 & 7,02 & 0,00 \\
     & Random & 76,08 & -0,47 & 1,42 & 8,07$\times$ & 1,56$\times$ & 40,1 & 40,2 & 6,51 & 7,42 & 0,00 \\
    \midrule
    \multirow{7}{*}{50\%} & Magnitude L1 & 71,49 & -5,06 & 0,98 & 11,72$\times$ & 3,96$\times$ & 50,1 & 79,7 & 5,43 & 6,04 & 0,00 \\
     & Magnitude L2 & 73,83 & -2,72 & 0,98 & 11,67$\times$ & 3,70$\times$ & 50,1 & 77,9 & 5,42 & 5,74 & 0,00 \\
     & BN-Scale & 74,62 & -1,93 & 1,13 & 10,17$\times$ & 2,86$\times$ & 50,1 & 69,8 & 5,43 & 6,33 & 0,00 \\
     & FPGM & 73,20 & -3,35 & 1,01 & 11,29$\times$ & 3,48$\times$ & 50,0 & 75,5 & 5,44 & 6,05 & 0,00 \\
     & Taylor & 76,15 & -0,40 & 1,29 & 8,90$\times$ & 1,94$\times$ & 50,1 & 53,4 & 5,45 & 6,21 & 0,00 \\
     & OBD-C & 75,54 & -1,01 & 1,24 & 9,24$\times$ & 2,00$\times$ & 50,2 & 54,4 & 5,43 & 6,09 & 0,00 \\
     & Random & 75,78 & -0,77 & 1,29 & 8,89$\times$ & 1,74$\times$ & 50,0 & 48,5 & 5,45 & 5,88 & 0,00 \\
    \midrule
    \multirow{7}{*}{60\%} & Magnitude L1 & 69,12 & -7,43 & 0,87 & 13,15$\times$ & 5,12$\times$ & 60,2 & 85,3 & 4,33 & 4,72 & 0,00 \\
     & Magnitude L2 & 69,48 & -7,07 & 0,87 & 13,10$\times$ & 4,96$\times$ & 60,1 & 84,8 & 4,35 & 4,83 & 0,00 \\
     & BN-Scale & 73,36 & -3,19 & 0,95 & 12,02$\times$ & 3,53$\times$ & 60,1 & 76,4 & 4,36 & 4,79 & 0,00 \\
     & FPGM & 71,49 & -5,06 & 0,87 & 13,17$\times$ & 4,56$\times$ & 60,0 & 82,8 & 4,35 & 4,57 & 0,00 \\
     & Taylor & 75,61 & -0,94 & 1,18 & 9,68$\times$ & 2,13$\times$ & 60,1 & 57,2 & 4,36 & 5,00 & 0,00 \\
     & OBD-C & 75,71 & -0,84 & 1,17 & 9,78$\times$ & 2,24$\times$ & 60,1 & 60,1 & 4,37 & 5,16 & 0,00 \\
     & Random & 75,63 & -0,92 & 1,20 & 9,56$\times$ & 2,16$\times$ & 60,1 & 59,2 & 4,37 & 5,16 & 0,00 \\
    \midrule
    \multirow{7}{*}{70\%} & Magnitude L1 & 67,88 & -8,67 & 0,77 & 14,81$\times$ & 6,55$\times$ & 70,1 & 88,9 & 3,27 & 3,48 & 0,00 \\
     & Magnitude L2 & 67,97 & -8,58 & 0,79 & 14,53$\times$ & 6,35$\times$ & 70,1 & 89,0 & 3,27 & 3,70 & 0,00 \\
     & BN-Scale & 71,56 & -4,99 & 0,86 & 13,27$\times$ & 5,02$\times$ & 70,1 & 84,1 & 3,28 & 3,86 & 0,00 \\
     & FPGM & 69,76 & -6,79 & 0,81 & 14,15$\times$ & 6,01$\times$ & 70,0 & 87,8 & 3,28 & 3,88 & 0,00 \\
     & Taylor & 75,07 & -1,48 & 1,08 & 10,60$\times$ & 2,42$\times$ & 70,1 & 63,9 & 3,29 & 3,83 & 0,00 \\
     & OBD-C & 74,96 & -1,59 & 1,06 & 10,80$\times$ & 2,65$\times$ & 70,1 & 66,7 & 3,28 & 3,79 & 0,00 \\
     & Random & 74,25 & -2,30 & 1,05 & 10,87$\times$ & 2,86$\times$ & 70,1 & 69,7 & 3,28 & 3,97 & 0,00 \\
    \midrule
    \multirow{7}{*}{80\%} & Magnitude L1 & 66,15 & -10,40 & 0,70 & 16,30$\times$ & 8,94$\times$ & 80,1 & 92,6 & 2,19 & 2,57 & 0,00 \\
     & Magnitude L2 & 66,30 & -10,25 & 0,70 & 16,42$\times$ & 8,81$\times$ & 80,0 & 92,6 & 2,20 & 2,53 & 0,00 \\
     & BN-Scale & 67,87 & -8,68 & 0,76 & 15,10$\times$ & 7,42$\times$ & 80,1 & 90,6 & 2,20 & 2,76 & 0,00 \\
     & FPGM & 66,89 & -9,66 & 0,69 & 16,55$\times$ & 8,67$\times$ & 80,1 & 92,5 & 2,19 & 2,53 & 0,00 \\
     & Taylor & 73,84 & -2,71 & 0,97 & 11,79$\times$ & 2,93$\times$ & 80,0 & 70,6 & 2,21 & 2,64 & 0,00 \\
     & OBD-C & 74,08 & -2,47 & 0,94 & 12,15$\times$ & 3,39$\times$ & 80,1 & 74,8 & 2,21 & 2,61 & 0,00 \\
     & Random & 73,30 & -3,25 & 0,88 & 12,95$\times$ & 3,81$\times$ & 80,1 & 78,6 & 2,20 & 2,57 & 0,00 \\
    \midrule
    \multirow{7}{*}{90\%} & Magnitude L1 & 62,24 & -14,31 & 0,62 & 18,35$\times$ & 14,87$\times$ & 90,1 & 96,3 & 1,10 & 1,42 & 0,00 \\
     & Magnitude L2 & 62,94 & -13,61 & 0,61 & 18,89$\times$ & 13,94$\times$ & 90,2 & 96,1 & 1,10 & 1,34 & 0,00 \\
     & BN-Scale & 63,46 & -13,09 & 0,63 & 18,23$\times$ & 12,23$\times$ & 90,1 & 95,3 & 1,11 & 1,30 & 0,00 \\
     & FPGM & 61,23 & -15,32 & 0,62 & 18,46$\times$ & 13,62$\times$ & 90,1 & 96,0 & 1,11 & 1,46 & 0,00 \\
     & Taylor & 72,84 & -3,71 & 0,84 & 13,65$\times$ & 4,04$\times$ & 90,1 & 80,6 & 1,12 & 1,47 & 0,00 \\
     & OBD-C & 72,26 & -4,29 & 0,78 & 14,60$\times$ & 4,33$\times$ & 90,0 & 81,9 & 1,13 & 1,40 & 0,00 \\
     & Random & 68,73 & -7,82 & 0,69 & 16,49$\times$ & 7,58$\times$ & 90,1 & 90,8 & 1,11 & 1,43 & 0,00 \\
    \bottomrule
  \end{tabular}}}
\end{table}

\clearpage
\renewcommand{\benchscale}{0.70}
\begin{table}[H]
  \centering
  \caption{Mesures complètes de ResNet-50 élagué, en INT8.}
  \label{tab:bench-resnet50}
  \centerline{\resizebox{\benchscale\linewidth}{!}{%
  \begin{tabular}{c l rr rrr rr rrr}
    \toprule
    & & \multicolumn{2}{c}{Précision} & \multicolumn{3}{c}{Temps} & \multicolumn{2}{c}{Réduction (\%)} & \multicolumn{3}{c}{Mémoire (Mio)} \\
    \cmidrule(lr){3-4} \cmidrule(lr){5-7} \cmidrule(lr){8-9} \cmidrule(lr){10-12}
    Cible & Critère & Top-1 & $\Delta$ Top-1 & Lat. (ms) & Accél. TPU & Accél. CPU & Param. & MACs & Taille & Interne & Externe \\
    \midrule
    -- & référence & 77,40 & +0,00 & 50,92 & 1,00$\times$ & 1,00$\times$ & 0,0 & 0,0 & 23,30 & 7,66 & 15,26 \\
    \midrule
    \multirow{7}{*}{10\%} & Magnitude L1 & 77,89 & +0,49 & 45,84 & 1,11$\times$ & 1,32$\times$ & 10,0 & 28,5 & 20,98 & 7,66 & 13,59 \\
     & Magnitude L2 & 77,64 & +0,24 & 45,15 & 1,13$\times$ & 1,31$\times$ & 10,0 & 27,1 & 20,97 & 7,66 & 13,50 \\
     & BN-Scale & 77,45 & +0,05 & 45,44 & 1,12$\times$ & 1,25$\times$ & 10,1 & 24,4 & 20,97 & 7,66 & 13,54 \\
     & FPGM & 77,38 & -0,02 & 43,90 & 1,16$\times$ & 1,24$\times$ & 10,0 & 20,2 & 20,96 & 7,65 & 13,02 \\
     & Taylor & 77,63 & +0,23 & 46,11 & 1,10$\times$ & 1,33$\times$ & 9,4 & 29,8 & 21,12 & 7,65 & 13,75 \\
     & OBD-C & 77,43 & +0,03 & 45,55 & 1,12$\times$ & 1,35$\times$ & 9,8 & 29,5 & 21,03 & 7,66 & 13,61 \\
     & Random & 77,57 & +0,17 & 47,00 & 1,08$\times$ & 1,07$\times$ & 10,1 & 10,4 & 20,99 & 7,65 & 14,04 \\
    \midrule
    \multirow{7}{*}{20\%} & Magnitude L1 & 77,66 & +0,26 & 39,00 & 1,31$\times$ & 1,55$\times$ & 20,0 & 40,5 & 18,65 & 7,66 & 11,57 \\
     & Magnitude L2 & 77,73 & +0,33 & 38,59 & 1,32$\times$ & 1,52$\times$ & 20,1 & 38,7 & 18,64 & 7,66 & 11,48 \\
     & BN-Scale & 77,49 & +0,09 & 38,41 & 1,33$\times$ & 1,50$\times$ & 20,0 & 37,1 & 18,65 & 7,66 & 11,45 \\
     & FPGM & 77,82 & +0,42 & 36,92 & 1,38$\times$ & 1,43$\times$ & 20,1 & 32,2 & 18,62 & 7,66 & 10,98 \\
     & Taylor & 77,47 & +0,07 & 38,76 & 1,31$\times$ & 1,62$\times$ & 20,0 & 43,1 & 18,66 & 7,66 & 11,54 \\
     & OBD-C & 77,60 & +0,20 & 37,94 & 1,34$\times$ & 1,60$\times$ & 20,1 & 41,6 & 18,64 & 7,66 & 11,29 \\
     & Random & 77,85 & +0,45 & 40,82 & 1,25$\times$ & 1,18$\times$ & 20,1 & 20,1 & 18,69 & 7,65 & 12,17 \\
    \midrule
    \multirow{7}{*}{30\%} & Magnitude L1 & 77,18 & -0,22 & 31,61 & 1,61$\times$ & 1,84$\times$ & 30,1 & 51,0 & 16,33 & 7,66 & 9,36 \\
     & Magnitude L2 & 77,40 & +0,00 & 31,26 & 1,63$\times$ & 1,80$\times$ & 30,1 & 49,7 & 16,32 & 7,66 & 9,27 \\
     & BN-Scale & 77,60 & +0,20 & 30,80 & 1,65$\times$ & 1,79$\times$ & 30,1 & 48,0 & 16,31 & 7,66 & 9,12 \\
     & FPGM & 77,76 & +0,36 & 29,40 & 1,73$\times$ & 1,74$\times$ & 30,0 & 45,3 & 16,31 & 7,66 & 8,69 \\
     & Taylor & 76,97 & -0,43 & 32,20 & 1,58$\times$ & 1,83$\times$ & 30,1 & 51,2 & 16,33 & 7,66 & 9,53 \\
     & OBD-C & 77,10 & -0,30 & 32,01 & 1,59$\times$ & 1,82$\times$ & 30,0 & 50,0 & 16,34 & 7,66 & 9,50 \\
     & Random & 76,89 & -0,51 & 33,13 & 1,54$\times$ & 1,33$\times$ & 30,1 & 30,4 & 16,38 & 7,66 & 9,73 \\
    \midrule
    \multirow{7}{*}{40\%} & Magnitude L1 & 77,31 & -0,09 & 23,66 & 2,15$\times$ & 2,26$\times$ & 40,1 & 61,2 & 14,00 & 7,66 & 6,95 \\
     & Magnitude L2 & 77,36 & -0,04 & 22,96 & 2,22$\times$ & 2,20$\times$ & 40,1 & 59,8 & 13,99 & 7,66 & 6,74 \\
     & BN-Scale & 77,47 & +0,07 & 23,50 & 2,17$\times$ & 2,18$\times$ & 40,1 & 59,1 & 13,99 & 7,66 & 6,89 \\
     & FPGM & 77,46 & +0,06 & 22,83 & 2,23$\times$ & 2,14$\times$ & 40,0 & 57,2 & 13,99 & 7,65 & 6,67 \\
     & Taylor & 77,20 & -0,20 & 24,27 & 2,10$\times$ & 2,06$\times$ & 40,2 & 57,1 & 14,03 & 7,66 & 7,11 \\
     & OBD-C & 77,16 & -0,24 & 24,66 & 2,07$\times$ & 2,05$\times$ & 40,1 & 56,2 & 14,03 & 7,66 & 7,25 \\
     & Random & 77,54 & +0,14 & 26,54 & 1,92$\times$ & 1,55$\times$ & 40,1 & 41,5 & 14,08 & 7,63 & 7,81 \\
    \midrule
    \multirow{7}{*}{50\%} & Magnitude L1 & 77,21 & -0,19 & 15,79 & 3,23$\times$ & 2,79$\times$ & 50,2 & 69,5 & 11,66 & 7,66 & 4,57 \\
     & Magnitude L2 & 77,15 & -0,25 & 16,52 & 3,08$\times$ & 2,75$\times$ & 50,0 & 68,9 & 11,70 & 7,69 & 4,77 \\
     & BN-Scale & 75,83 & -1,57 & 15,90 & 3,20$\times$ & 2,73$\times$ & 50,1 & 68,7 & 11,65 & 7,66 & 4,62 \\
     & FPGM & 77,14 & -0,26 & 15,03 & 3,39$\times$ & 2,79$\times$ & 50,0 & 68,3 & 11,66 & 7,66 & 4,34 \\
     & Taylor & 77,38 & -0,02 & 16,20 & 3,14$\times$ & 2,30$\times$ & 50,2 & 62,3 & 11,74 & 7,66 & 4,63 \\
     & OBD-C & 77,04 & -0,36 & 16,00 & 3,18$\times$ & 2,30$\times$ & 50,2 & 62,1 & 11,71 & 7,66 & 4,51 \\
     & Random & 76,87 & -0,53 & 16,12 & 3,16$\times$ & 1,77$\times$ & 50,2 & 50,0 & 11,77 & 7,65 & 4,52 \\
    \midrule
    \multirow{7}{*}{60\%} & Magnitude L1 & 77,12 & -0,28 & 7,96 & 6,40$\times$ & 3,78$\times$ & 60,2 & 78,2 & 9,34 & 7,73 & 2,03 \\
     & Magnitude L2 & 77,05 & -0,35 & 8,56 & 5,95$\times$ & 3,71$\times$ & 60,1 & 77,8 & 9,35 & 7,70 & 2,25 \\
     & BN-Scale & 75,88 & -1,52 & 8,80 & 5,79$\times$ & 3,68$\times$ & 60,2 & 77,4 & 9,31 & 7,66 & 2,41 \\
     & FPGM & 76,51 & -0,89 & 8,67 & 5,88$\times$ & 3,90$\times$ & 60,1 & 78,4 & 9,34 & 7,70 & 2,34 \\
     & Taylor & 77,44 & +0,04 & 9,22 & 5,52$\times$ & 2,56$\times$ & 60,2 & 67,0 & 9,44 & 7,66 & 2,40 \\
     & OBD-C & 77,31 & -0,09 & 8,85 & 5,75$\times$ & 2,63$\times$ & 60,1 & 67,4 & 9,44 & 7,66 & 2,35 \\
     & Random & 76,57 & -0,83 & 10,85 & 4,69$\times$ & 2,16$\times$ & 60,1 & 61,2 & 9,47 & 7,66 & 2,82 \\
    \midrule
    \multirow{7}{*}{70\%} & Magnitude L1 & 75,54 & -1,86 & 1,49 & 34,12$\times$ & 5,50$\times$ & 70,1 & 85,7 & 7,02 & 7,68 & 0,00 \\
     & Magnitude L2 & 75,43 & -1,97 & 1,50 & 34,02$\times$ & 5,55$\times$ & 70,1 & 86,1 & 7,02 & 7,71 & 0,00 \\
     & BN-Scale & 75,67 & -1,73 & 1,57 & 32,42$\times$ & 5,35$\times$ & 70,1 & 85,2 & 7,02 & 7,73 & 0,03 \\
     & FPGM & 75,60 & -1,80 & 1,37 & 37,16$\times$ & 6,01$\times$ & 70,2 & 87,1 & 6,98 & 7,52 & 0,00 \\
     & Taylor & 76,78 & -0,62 & 2,59 & 19,66$\times$ & 2,96$\times$ & 70,0 & 72,4 & 7,18 & 7,66 & 0,23 \\
     & OBD-C & 76,87 & -0,53 & 2,18 & 23,37$\times$ & 3,10$\times$ & 70,0 & 73,0 & 7,15 & 7,66 & 0,05 \\
     & Random & 76,56 & -0,84 & 2,94 & 17,29$\times$ & 2,74$\times$ & 70,0 & 70,6 & 7,17 & 7,66 & 0,43 \\
    \midrule
    \multirow{7}{*}{80\%} & Magnitude L1 & 73,91 & -3,49 & 1,23 & 41,52$\times$ & 8,37$\times$ & 80,2 & 91,9 & 4,68 & 5,15 & 0,00 \\
     & Magnitude L2 & 73,89 & -3,51 & 1,15 & 44,41$\times$ & 9,02$\times$ & 80,2 & 92,5 & 4,66 & 5,18 & 0,00 \\
     & BN-Scale & 74,92 & -2,48 & 1,22 & 41,74$\times$ & 8,18$\times$ & 80,1 & 91,5 & 4,68 & 5,26 & 0,00 \\
     & FPGM & 73,03 & -4,37 & 1,08 & 47,31$\times$ & 9,49$\times$ & 80,2 & 93,2 & 4,63 & 5,05 & 0,00 \\
     & Taylor & 76,08 & -1,32 & 1,92 & 26,54$\times$ & 3,69$\times$ & 80,1 & 79,4 & 4,85 & 5,57 & 0,00 \\
     & OBD-C & 76,20 & -1,20 & 1,84 & 27,67$\times$ & 3,89$\times$ & 80,1 & 80,2 & 4,81 & 5,38 & 0,00 \\
     & Random & 76,23 & -1,17 & 1,78 & 28,56$\times$ & 3,65$\times$ & 80,2 & 80,1 & 4,81 & 5,17 & 0,00 \\
    \midrule
    \multirow{7}{*}{90\%} & Magnitude L1 & 70,24 & -7,16 & 0,95 & 53,71$\times$ & 14,52$\times$ & 90,1 & 96,4 & 2,36 & 2,66 & 0,00 \\
     & Magnitude L2 & 71,41 & -5,99 & 0,93 & 54,84$\times$ & 14,67$\times$ & 90,1 & 96,3 & 2,38 & 2,69 & 0,00 \\
     & BN-Scale & 68,81 & -8,59 & 0,96 & 53,13$\times$ & 14,09$\times$ & 90,2 & 96,1 & 2,36 & 2,78 & 0,00 \\
     & FPGM & 70,94 & -6,46 & 0,94 & 54,27$\times$ & 14,13$\times$ & 90,2 & 96,3 & 2,35 & 2,96 & 0,00 \\
     & Taylor & 74,30 & -3,10 & 1,52 & 33,54$\times$ & 5,41$\times$ & 90,1 & 88,1 & 2,50 & 2,96 & 0,00 \\
     & OBD-C & 75,17 & -2,23 & 1,36 & 37,37$\times$ & 6,03$\times$ & 90,2 & 89,2 & 2,46 & 2,92 & 0,00 \\
     & Random & 73,69 & -3,71 & 1,34 & 38,05$\times$ & 6,30$\times$ & 90,1 & 90,6 & 2,48 & 3,02 & 0,00 \\
    \bottomrule
  \end{tabular}}}
\end{table}

\clearpage
\renewcommand{\benchscale}{0.78}
\begin{table}[H]
  \centering
  \caption{Mesures complètes de VGG-19 élagué, en INT8.}
  \label{tab:bench-vgg19}
  \centerline{\resizebox{\benchscale\linewidth}{!}{%
  \begin{tabular}{c l rr rrr rr rrr}
    \toprule
    & & \multicolumn{2}{c}{Précision} & \multicolumn{3}{c}{Temps} & \multicolumn{2}{c}{Réduction (\%)} & \multicolumn{3}{c}{Mémoire (Mio)} \\
    \cmidrule(lr){3-4} \cmidrule(lr){5-7} \cmidrule(lr){8-9} \cmidrule(lr){10-12}
    Cible & Critère & Top-1 & $\Delta$ Top-1 & Lat. (ms) & Accél. TPU & Accél. CPU & Param. & MACs & Taille & Interne & Externe \\
    \midrule
    -- & référence & 73,13 & +0,00 & 39,29 & 1,00$\times$ & 1,00$\times$ & 0,0 & 0,0 & 19,30 & 7,58 & 11,64 \\
    \midrule
    \multirow{6}{*}{10\%} & Magnitude L1 & 74,00 & +0,87 & 34,92 & 1,13$\times$ & 1,11$\times$ & 10,1 & 11,4 & 17,37 & 7,57 & 10,42 \\
     & Magnitude L2 & 73,55 & +0,42 & 35,02 & 1,12$\times$ & 1,12$\times$ & 10,1 & 12,2 & 17,36 & 7,63 & 10,35 \\
     & BN-Scale & 73,57 & +0,44 & 35,41 & 1,11$\times$ & 1,10$\times$ & 10,0 & 11,1 & 17,37 & 7,54 & 10,45 \\
     & FPGM & 73,77 & +0,64 & 35,25 & 1,11$\times$ & 1,13$\times$ & 10,1 & 12,1 & 17,36 & 7,59 & 10,41 \\
     & Taylor & 73,59 & +0,46 & 35,40 & 1,11$\times$ & 1,12$\times$ & 10,1 & 11,3 & 17,36 & 7,59 & 10,45 \\
     & Random & 73,89 & +0,76 & 36,34 & 1,08$\times$ & 1,08$\times$ & 10,0 & 9,9 & 17,37 & 7,56 & 10,75 \\
    \midrule
    \multirow{6}{*}{20\%} & Magnitude L1 & 73,58 & +0,45 & 29,13 & 1,35$\times$ & 1,22$\times$ & 20,1 & 19,0 & 15,44 & 7,72 & 8,55 \\
     & Magnitude L2 & 73,74 & +0,61 & 28,99 & 1,36$\times$ & 1,23$\times$ & 20,1 & 18,9 & 15,44 & 7,72 & 8,50 \\
     & BN-Scale & 73,62 & +0,49 & 28,62 & 1,37$\times$ & 1,22$\times$ & 20,0 & 18,1 & 15,45 & 7,59 & 8,39 \\
     & FPGM & 73,29 & +0,16 & 29,75 & 1,32$\times$ & 1,22$\times$ & 20,1 & 18,7 & 15,44 & 7,59 & 8,73 \\
     & Taylor & 74,01 & +0,88 & 29,52 & 1,33$\times$ & 1,24$\times$ & 20,1 & 20,2 & 15,44 & 7,72 & 8,66 \\
     & Random & 73,18 & +0,05 & 32,94 & 1,19$\times$ & 1,20$\times$ & 20,0 & 19,8 & 15,45 & 7,57 & 9,70 \\
    \midrule
    \multirow{6}{*}{30\%} & Magnitude L1 & 74,02 & +0,89 & 22,19 & 1,77$\times$ & 1,33$\times$ & 30,0 & 25,6 & 13,53 & 7,70 & 6,45 \\
     & Magnitude L2 & 73,87 & +0,74 & 22,12 & 1,78$\times$ & 1,33$\times$ & 30,0 & 25,3 & 13,53 & 7,68 & 6,43 \\
     & BN-Scale & 73,24 & +0,11 & 23,26 & 1,69$\times$ & 1,30$\times$ & 30,1 & 24,6 & 13,51 & 7,59 & 6,76 \\
     & FPGM & 73,36 & +0,23 & 22,30 & 1,76$\times$ & 1,32$\times$ & 30,0 & 25,3 & 13,53 & 7,68 & 6,48 \\
     & Taylor & 73,52 & +0,39 & 23,27 & 1,69$\times$ & 1,38$\times$ & 30,1 & 28,9 & 13,51 & 7,73 & 6,77 \\
     & Random & 73,21 & +0,08 & 23,49 & 1,67$\times$ & 1,37$\times$ & 30,0 & 29,2 & 13,53 & 7,56 & 6,83 \\
    \midrule
    \multirow{6}{*}{40\%} & Magnitude L1 & 73,74 & +0,61 & 16,34 & 2,40$\times$ & 1,45$\times$ & 40,1 & 32,4 & 11,60 & 7,71 & 4,66 \\
     & Magnitude L2 & 73,78 & +0,65 & 16,29 & 2,41$\times$ & 1,45$\times$ & 40,1 & 32,1 & 11,60 & 7,72 & 4,65 \\
     & BN-Scale & 73,99 & +0,86 & 16,71 & 2,35$\times$ & 1,45$\times$ & 40,1 & 31,9 & 11,60 & 7,71 & 4,78 \\
     & FPGM & 73,33 & +0,20 & 15,66 & 2,51$\times$ & 1,47$\times$ & 40,1 & 32,2 & 11,59 & 7,68 & 4,48 \\
     & Taylor & 73,64 & +0,51 & 16,41 & 2,39$\times$ & 1,56$\times$ & 40,1 & 37,7 & 11,59 & 7,67 & 4,68 \\
     & Random & 72,79 & -0,34 & 19,06 & 2,06$\times$ & 1,63$\times$ & 40,1 & 40,5 & 11,58 & 7,57 & 5,50 \\
    \midrule
    \multirow{6}{*}{50\%} & Magnitude L1 & 73,88 & +0,75 & 10,34 & 3,80$\times$ & 1,63$\times$ & 50,0 & 38,8 & 9,68 & 7,71 & 2,79 \\
     & Magnitude L2 & 73,11 & -0,02 & 9,51 & 4,13$\times$ & 1,64$\times$ & 50,1 & 39,3 & 9,67 & 7,68 & 2,55 \\
     & BN-Scale & 72,91 & -0,22 & 9,26 & 4,24$\times$ & 1,64$\times$ & 50,2 & 39,1 & 9,65 & 7,71 & 2,45 \\
     & FPGM & 73,53 & +0,40 & 10,08 & 3,90$\times$ & 1,66$\times$ & 50,1 & 39,4 & 9,66 & 7,71 & 2,70 \\
     & Taylor & 72,55 & -0,58 & 11,54 & 3,40$\times$ & 1,82$\times$ & 50,0 & 47,1 & 9,68 & 7,72 & 3,18 \\
     & Random & 72,05 & -1,08 & 10,21 & 3,85$\times$ & 1,88$\times$ & 50,1 & 50,7 & 9,67 & 7,57 & 2,75 \\
    \midrule
    \multirow{6}{*}{60\%} & Magnitude L1 & 73,66 & +0,53 & 2,07 & 18,98$\times$ & 1,89$\times$ & 60,1 & 46,1 & 7,74 & 7,67 & 0,33 \\
     & Magnitude L2 & 73,37 & +0,24 & 2,93 & 13,40$\times$ & 1,89$\times$ & 60,1 & 46,5 & 7,74 & 7,68 & 0,55 \\
     & BN-Scale & 73,06 & -0,07 & 2,91 & 13,51$\times$ & 1,86$\times$ & 60,1 & 46,1 & 7,75 & 7,73 & 0,56 \\
     & FPGM & 73,09 & -0,04 & 2,73 & 14,39$\times$ & 1,87$\times$ & 60,1 & 46,4 & 7,74 & 7,72 & 0,47 \\
     & Taylor & 72,70 & -0,43 & 2,65 & 14,85$\times$ & 2,18$\times$ & 60,1 & 56,4 & 7,74 & 7,73 & 0,48 \\
     & Random & 71,17 & -1,96 & 4,97 & 7,90$\times$ & 2,31$\times$ & 60,2 & 60,2 & 7,73 & 7,58 & 1,08 \\
    \midrule
    \multirow{6}{*}{70\%} & Magnitude L1 & 72,84 & -0,29 & 1,04 & 37,70$\times$ & 2,20$\times$ & 70,1 & 54,1 & 5,82 & 6,26 & 0,00 \\
     & Magnitude L2 & 73,00 & -0,13 & 1,04 & 37,88$\times$ & 2,24$\times$ & 70,1 & 54,2 & 5,81 & 6,30 & 0,00 \\
     & BN-Scale & 72,63 & -0,50 & 1,04 & 37,80$\times$ & 2,25$\times$ & 70,1 & 55,0 & 5,82 & 6,24 & 0,00 \\
     & FPGM & 72,76 & -0,37 & 1,01 & 38,91$\times$ & 2,26$\times$ & 70,1 & 54,3 & 5,81 & 6,14 & 0,00 \\
     & Taylor & 71,83 & -1,30 & 1,00 & 39,44$\times$ & 2,79$\times$ & 70,0 & 66,2 & 5,82 & 6,81 & 0,00 \\
     & Random & 71,01 & -2,12 & 1,02 & 38,37$\times$ & 3,06$\times$ & 70,2 & 69,9 & 5,80 & 6,83 & 0,00 \\
    \midrule
    \multirow{6}{*}{80\%} & Magnitude L1 & 72,40 & -0,73 & 0,88 & 44,46$\times$ & 2,93$\times$ & 80,0 & 65,2 & 3,89 & 4,24 & 0,00 \\
     & Magnitude L2 & 72,29 & -0,84 & 0,86 & 45,93$\times$ & 3,00$\times$ & 80,0 & 66,1 & 3,89 & 4,19 & 0,00 \\
     & BN-Scale & 72,41 & -0,72 & 0,87 & 45,08$\times$ & 3,04$\times$ & 80,1 & 67,0 & 3,89 & 4,26 & 0,00 \\
     & FPGM & 72,37 & -0,76 & 0,84 & 46,80$\times$ & 3,16$\times$ & 80,2 & 68,2 & 3,86 & 4,12 & 0,00 \\
     & Taylor & 70,05 & -3,08 & 0,82 & 47,84$\times$ & 3,95$\times$ & 80,2 & 76,9 & 3,87 & 4,50 & 0,00 \\
     & Random & 69,77 & -3,36 & 0,80 & 49,05$\times$ & 4,47$\times$ & 80,0 & 80,1 & 3,90 & 4,38 & 0,00 \\
    \midrule
    \multirow{6}{*}{90\%} & Magnitude L1 & 70,79 & -2,34 & 0,67 & 58,52$\times$ & 5,41$\times$ & 90,1 & 82,0 & 1,94 & 2,27 & 0,00 \\
     & Magnitude L2 & 69,70 & -3,43 & 0,68 & 58,08$\times$ & 5,39$\times$ & 90,1 & 81,8 & 1,94 & 2,30 & 0,00 \\
     & BN-Scale & 68,35 & -4,78 & 0,67 & 58,92$\times$ & 5,76$\times$ & 90,1 & 83,1 & 1,94 & 2,36 & 0,00 \\
     & FPGM & 60,11 & -13,02 & 0,61 & 64,31$\times$ & 8,41$\times$ & 90,1 & 90,1 & 1,95 & 2,12 & 0,00 \\
     & Taylor & 65,82 & -7,31 & 0,65 & 60,61$\times$ & 6,63$\times$ & 90,1 & 87,1 & 1,95 & 2,35 & 0,00 \\
     & Random & 63,53 & -9,60 & 0,64 & 61,37$\times$ & 8,21$\times$ & 90,1 & 90,2 & 1,95 & 2,42 & 0,00 \\
    \bottomrule
  \end{tabular}}}
\end{table}

\clearpage
\renewcommand{\benchscale}{0.72}
\begin{table}[H]
  \centering
  \caption{Mesures complètes de WRN-28-10 élagué, en INT8.}
  \label{tab:bench-wrn_28_10}
  \centerline{\resizebox{\benchscale\linewidth}{!}{%
  \begin{tabular}{c l rr rrr rr rrr}
    \toprule
    & & \multicolumn{2}{c}{Précision} & \multicolumn{3}{c}{Temps} & \multicolumn{2}{c}{Réduction (\%)} & \multicolumn{3}{c}{Mémoire (Mio)} \\
    \cmidrule(lr){3-4} \cmidrule(lr){5-7} \cmidrule(lr){8-9} \cmidrule(lr){10-12}
    Cible & Critère & Top-1 & $\Delta$ Top-1 & Lat. (ms) & Accél. TPU & Accél. CPU & Param. & MACs & Taille & Interne & Externe \\
    \midrule
    -- & référence & 80,45 & +0,00 & 92,13 & 1,00$\times$ & 1,00$\times$ & 0,0 & 0,0 & 35,13 & 7,65 & 27,74 \\
    \midrule
    \multirow{7}{*}{10\%} & Magnitude L1 & 80,76 & +0,31 & 82,98 & 1,11$\times$ & 1,12$\times$ & 10,0 & 11,8 & 31,63 & 7,65 & 24,80 \\
     & Magnitude L2 & 80,75 & +0,30 & 81,88 & 1,13$\times$ & 1,11$\times$ & 10,0 & 11,4 & 31,63 & 7,63 & 24,64 \\
     & BN-Scale & 80,74 & +0,29 & 83,42 & 1,10$\times$ & 1,13$\times$ & 10,0 & 12,9 & 31,63 & 7,65 & 24,96 \\
     & FPGM & 80,77 & +0,32 & 81,30 & 1,13$\times$ & 1,05$\times$ & 10,0 & 4,5 & 31,63 & 7,65 & 24,35 \\
     & Taylor & 80,76 & +0,31 & 60,44 & 1,52$\times$ & 1,30$\times$ & 28,8 & 24,2 & 25,06 & 7,65 & 17,96 \\
     & OBD-C & 80,60 & +0,15 & 82,42 & 1,12$\times$ & 1,12$\times$ & 9,9 & 12,2 & 31,66 & 7,65 & 24,65 \\
     & Random & 80,42 & -0,03 & 86,61 & 1,06$\times$ & 1,07$\times$ & 10,0 & 9,5 & 31,63 & 7,65 & 26,05 \\
    \midrule
    \multirow{5}{*}{20\%} & Magnitude L1 & 80,50 & +0,05 & 70,76 & 1,30$\times$ & 1,22$\times$ & 20,0 & 19,0 & 28,13 & 7,65 & 21,28 \\
     & Magnitude L2 & 80,53 & +0,08 & 71,14 & 1,29$\times$ & 1,22$\times$ & 20,0 & 18,9 & 28,13 & 7,65 & 21,29 \\
     & FPGM & 80,64 & +0,19 & 73,35 & 1,26$\times$ & 1,17$\times$ & 20,0 & 15,8 & 28,13 & 7,63 & 21,92 \\
     & OBD-C & 81,00 & +0,55 & 71,31 & 1,29$\times$ & 1,22$\times$ & 20,0 & 19,2 & 28,13 & 7,65 & 21,28 \\
     & Random & 80,50 & +0,05 & 73,57 & 1,25$\times$ & 1,22$\times$ & 20,0 & 20,1 & 28,13 & 7,64 & 22,07 \\
    \midrule
    \multirow{7}{*}{30\%} & Magnitude L1 & 80,41 & -0,04 & 60,41 & 1,53$\times$ & 1,31$\times$ & 30,1 & 25,8 & 24,60 & 7,62 & 17,92 \\
     & Magnitude L2 & 80,56 & +0,11 & 60,45 & 1,52$\times$ & 1,32$\times$ & 30,1 & 25,8 & 24,63 & 7,65 & 17,93 \\
     & BN-Scale & 80,72 & +0,27 & 60,33 & 1,53$\times$ & 1,30$\times$ & 30,1 & 25,2 & 24,62 & 7,62 & 17,91 \\
     & FPGM & 80,41 & -0,04 & 59,99 & 1,54$\times$ & 1,26$\times$ & 30,1 & 20,4 & 24,60 & 7,65 & 17,73 \\
     & Taylor & 80,79 & +0,34 & 60,05 & 1,53$\times$ & 1,31$\times$ & 30,1 & 25,6 & 24,59 & 7,65 & 17,85 \\
     & OBD-C & 81,14 & +0,69 & 60,31 & 1,53$\times$ & 1,32$\times$ & 30,0 & 26,7 & 24,63 & 7,65 & 17,89 \\
     & Random & 80,42 & -0,03 & 66,36 & 1,39$\times$ & 1,39$\times$ & 30,1 & 30,1 & 24,60 & 7,63 & 19,80 \\
    \midrule
    \multirow{7}{*}{40\%} & Magnitude L1 & 80,38 & -0,07 & 49,28 & 1,87$\times$ & 1,72$\times$ & 40,2 & 44,3 & 21,07 & 7,65 & 14,51 \\
     & Magnitude L2 & 80,55 & +0,10 & 49,48 & 1,86$\times$ & 1,54$\times$ & 40,0 & 37,3 & 21,13 & 7,65 & 14,63 \\
     & BN-Scale & 80,45 & +0,00 & 48,37 & 1,90$\times$ & 1,42$\times$ & 40,3 & 30,6 & 21,04 & 7,65 & 14,29 \\
     & FPGM & 80,57 & +0,12 & 49,48 & 1,86$\times$ & 1,35$\times$ & 40,2 & 25,9 & 21,08 & 7,64 & 14,61 \\
     & Taylor & 80,36 & -0,09 & 51,77 & 1,78$\times$ & 1,55$\times$ & 40,1 & 38,0 & 21,10 & 7,65 & 15,24 \\
     & OBD-C & 80,32 & -0,13 & 48,04 & 1,92$\times$ & 1,58$\times$ & 40,1 & 40,4 & 21,12 & 7,65 & 14,27 \\
     & Random & 79,79 & -0,66 & 48,33 & 1,91$\times$ & 1,66$\times$ & 40,0 & 40,5 & 21,12 & 7,58 & 14,30 \\
    \midrule
    \multirow{7}{*}{50\%} & Magnitude L1 & 78,52 & -1,93 & 40,19 & 2,29$\times$ & 2,26$\times$ & 50,0 & 58,6 & 17,63 & 7,65 & 11,77 \\
     & Magnitude L2 & 80,18 & -0,27 & 38,12 & 2,42$\times$ & 1,68$\times$ & 50,1 & 43,3 & 17,61 & 7,65 & 11,12 \\
     & BN-Scale & 80,51 & +0,06 & 37,25 & 2,47$\times$ & 1,51$\times$ & 50,2 & 34,9 & 17,59 & 7,63 & 10,77 \\
     & FPGM & 80,61 & +0,16 & 35,63 & 2,59$\times$ & 1,48$\times$ & 50,1 & 32,8 & 17,61 & 7,64 & 10,36 \\
     & Taylor & 80,40 & -0,05 & 38,99 & 2,36$\times$ & 1,81$\times$ & 50,2 & 46,1 & 17,57 & 7,65 & 11,39 \\
     & OBD-C & 80,16 & -0,29 & 37,69 & 2,44$\times$ & 1,87$\times$ & 50,2 & 47,9 & 17,57 & 7,65 & 11,01 \\
     & Random & 80,09 & -0,36 & 38,21 & 2,41$\times$ & 1,97$\times$ & 50,1 & 48,8 & 17,60 & 7,65 & 11,27 \\
    \midrule
    \multirow{7}{*}{60\%} & Magnitude L1 & 73,95 & -6,50 & 27,07 & 3,40$\times$ & 3,24$\times$ & 60,1 & 69,9 & 14,09 & 7,65 & 7,75 \\
     & Magnitude L2 & 80,17 & -0,28 & 26,54 & 3,47$\times$ & 1,88$\times$ & 60,1 & 48,8 & 14,12 & 7,65 & 7,52 \\
     & BN-Scale & 80,02 & -0,43 & 28,01 & 3,29$\times$ & 1,75$\times$ & 60,2 & 44,2 & 14,09 & 7,62 & 8,05 \\
     & FPGM & 79,84 & -0,61 & 25,60 & 3,60$\times$ & 1,66$\times$ & 60,0 & 41,2 & 14,12 & 7,63 & 7,25 \\
     & Taylor & 80,29 & -0,16 & 26,37 & 3,49$\times$ & 2,03$\times$ & 60,2 & 52,9 & 14,07 & 7,64 & 7,50 \\
     & OBD-C & 80,11 & -0,34 & 28,52 & 3,23$\times$ & 2,11$\times$ & 60,0 & 54,6 & 14,13 & 7,65 & 8,19 \\
     & Random & 79,46 & -0,99 & 29,40 & 3,13$\times$ & 2,45$\times$ & 60,2 & 59,8 & 14,08 & 7,65 & 8,50 \\
    \midrule
    \multirow{7}{*}{70\%} & Magnitude L1 & 72,68 & -7,77 & 14,80 & 6,22$\times$ & 4,90$\times$ & 70,1 & 80,9 & 10,58 & 7,65 & 4,11 \\
     & Magnitude L2 & 79,05 & -1,40 & 15,18 & 6,07$\times$ & 2,51$\times$ & 70,3 & 61,7 & 10,53 & 7,64 & 4,05 \\
     & BN-Scale & 79,63 & -0,82 & 13,47 & 6,84$\times$ & 2,51$\times$ & 70,2 & 61,7 & 10,55 & 7,65 & 3,59 \\
     & FPGM & 79,23 & -1,22 & 16,12 & 5,71$\times$ & 2,26$\times$ & 70,0 & 58,3 & 10,63 & 7,65 & 4,19 \\
     & Taylor & 80,36 & -0,09 & 15,02 & 6,13$\times$ & 2,31$\times$ & 70,0 & 58,6 & 10,64 & 7,64 & 3,96 \\
     & OBD-C & 80,25 & -0,20 & 14,20 & 6,49$\times$ & 2,54$\times$ & 70,2 & 62,0 & 10,55 & 7,65 & 3,82 \\
     & Random & 79,13 & -1,32 & 15,25 & 6,04$\times$ & 3,38$\times$ & 70,0 & 70,9 & 10,62 & 7,65 & 4,20 \\
    \midrule
    \multirow{7}{*}{80\%} & Magnitude L1 & 71,17 & -9,28 & 2,23 & 41,37$\times$ & 8,16$\times$ & 80,2 & 89,1 & 7,04 & 7,66 & 0,25 \\
     & Magnitude L2 & 78,67 & -1,78 & 3,19 & 28,84$\times$ & 4,20$\times$ & 80,2 & 77,1 & 7,06 & 7,65 & 0,13 \\
     & BN-Scale & 76,86 & -3,59 & 3,63 & 25,41$\times$ & 5,66$\times$ & 80,1 & 83,9 & 7,07 & 7,66 & 0,54 \\
     & FPGM & 78,33 & -2,12 & 3,20 & 28,75$\times$ & 3,61$\times$ & 80,3 & 74,3 & 7,01 & 7,65 & 0,07 \\
     & Taylor & 79,25 & -1,20 & 4,92 & 18,74$\times$ & 2,94$\times$ & 80,1 & 67,6 & 7,10 & 7,65 & 0,62 \\
     & OBD-C & 79,35 & -1,10 & 4,37 & 21,08$\times$ & 3,11$\times$ & 80,1 & 69,3 & 7,10 & 7,65 & 0,35 \\
     & Random & 77,87 & -2,58 & 4,36 & 21,14$\times$ & 4,71$\times$ & 80,1 & 80,0 & 7,07 & 7,65 & 0,71 \\
    \midrule
    \multirow{7}{*}{90\%} & Magnitude L1 & 67,58 & -12,87 & 1,40 & 65,96$\times$ & 15,24$\times$ & 90,2 & 94,8 & 3,53 & 4,21 & 0,03 \\
     & Magnitude L2 & 72,01 & -8,44 & 1,56 & 59,02$\times$ & 13,19$\times$ & 90,3 & 93,7 & 3,49 & 4,38 & 0,02 \\
     & BN-Scale & 72,66 & -7,79 & 1,53 & 60,13$\times$ & 11,90$\times$ & 90,0 & 93,3 & 3,59 & 4,22 & 0,03 \\
     & FPGM & 75,63 & -4,82 & 1,84 & 50,06$\times$ & 8,97$\times$ & 90,1 & 89,6 & 3,57 & 4,28 & 0,03 \\
     & Taylor & 78,10 & -2,35 & 2,93 & 31,42$\times$ & 4,37$\times$ & 90,0 & 78,5 & 3,60 & 4,39 & 0,06 \\
     & OBD-C & 77,48 & -2,97 & 2,80 & 32,91$\times$ & 4,96$\times$ & 90,1 & 81,3 & 3,57 & 4,31 & 0,05 \\
     & Random & 75,96 & -4,49 & 1,92 & 48,02$\times$ & 8,61$\times$ & 90,1 & 89,6 & 3,56 & 4,03 & 0,04 \\
    \bottomrule
  \end{tabular}}}
\end{table}

\clearpage
\renewcommand{\benchscale}{0.82}
\begin{table}[H]
  \centering
  \caption{Mesures complètes de MobileNetV2 élagué, en INT8.}
  \label{tab:bench-mobilenetv2}
  \centerline{\resizebox{\benchscale\linewidth}{!}{%
  \begin{tabular}{c l rr rrr rr rrr}
    \toprule
    & & \multicolumn{2}{c}{Précision} & \multicolumn{3}{c}{Temps} & \multicolumn{2}{c}{Réduction (\%)} & \multicolumn{3}{c}{Mémoire (Mio)} \\
    \cmidrule(lr){3-4} \cmidrule(lr){5-7} \cmidrule(lr){8-9} \cmidrule(lr){10-12}
    Cible & Critère & Top-1 & $\Delta$ Top-1 & Lat. (ms) & Accél. TPU & Accél. CPU & Param. & MACs & Taille & Interne & Externe \\
    \midrule
    -- & référence & 62,38 & +0,00 & 1,24 & 1,00$\times$ & 1,00$\times$ & 0,0 & 0,0 & 2,70 & 2,79 & 0,00 \\
    \midrule
    \multirow{6}{*}{10\%} & Magnitude L1 & 64,47 & +2,09 & 1,16 & 1,07$\times$ & 1,18$\times$ & 10,0 & 16,6 & 2,43 & 2,53 & 0,00 \\
     & Magnitude L2 & 64,41 & +2,03 & 1,17 & 1,06$\times$ & 1,16$\times$ & 9,7 & 16,4 & 2,44 & 2,54 & 0,00 \\
     & BN-Scale & 64,21 & +1,83 & 1,16 & 1,07$\times$ & 1,18$\times$ & 6,8 & 22,2 & 2,50 & 2,60 & 0,00 \\
     & FPGM & 64,40 & +2,02 & 1,18 & 1,05$\times$ & 1,03$\times$ & 13,6 & 5,1 & 2,35 & 2,41 & 0,00 \\
     & Taylor & 63,83 & +1,45 & 1,00 & 1,24$\times$ & 1,32$\times$ & 34,7 & 37,9 & 1,76 & 1,84 & 0,00 \\
     & Random & 64,44 & +2,06 & 1,25 & 0,99$\times$ & 0,99$\times$ & 10,1 & 10,5 & 2,45 & 2,64 & 0,00 \\
    \midrule
    \multirow{5}{*}{20\%} & Magnitude L1 & 64,06 & +1,68 & 1,09 & 1,14$\times$ & 1,19$\times$ & 20,0 & 26,3 & 2,16 & 2,26 & 0,00 \\
     & Magnitude L2 & 64,78 & +2,40 & 1,13 & 1,10$\times$ & 1,25$\times$ & 19,7 & 25,4 & 2,17 & 2,26 & 0,00 \\
     & BN-Scale & 65,14 & +2,76 & 1,09 & 1,14$\times$ & 1,30$\times$ & 20,0 & 29,1 & 2,17 & 2,26 & 0,00 \\
     & FPGM & 64,16 & +1,78 & 1,12 & 1,11$\times$ & 1,07$\times$ & 25,9 & 9,8 & 2,03 & 2,07 & 0,00 \\
     & Random & 64,00 & +1,62 & 1,20 & 1,04$\times$ & 1,04$\times$ & 20,1 & 19,6 & 2,20 & 2,43 & 0,00 \\
    \midrule
    \multirow{5}{*}{30\%} & Magnitude L1 & 64,99 & +2,61 & 1,04 & 1,19$\times$ & 1,28$\times$ & 30,0 & 35,0 & 1,89 & 1,97 & 0,00 \\
     & Magnitude L2 & 64,68 & +2,30 & 1,04 & 1,19$\times$ & 1,27$\times$ & 30,0 & 33,3 & 1,89 & 1,99 & 0,00 \\
     & BN-Scale & 65,21 & +2,83 & 1,02 & 1,22$\times$ & 1,41$\times$ & 30,0 & 39,2 & 1,89 & 1,98 & 0,00 \\
     & FPGM & 64,16 & +1,78 & 1,09 & 1,13$\times$ & 1,10$\times$ & 30,4 & 16,6 & 1,90 & 1,94 & 0,00 \\
     & Random & 63,35 & +0,97 & 1,16 & 1,07$\times$ & 1,10$\times$ & 30,2 & 29,1 & 1,95 & 2,22 & 0,00 \\
    \midrule
    \multirow{6}{*}{40\%} & Magnitude L1 & 64,05 & +1,67 & 0,97 & 1,28$\times$ & 1,39$\times$ & 40,1 & 43,6 & 1,62 & 1,70 & 0,00 \\
     & Magnitude L2 & 64,20 & +1,82 & 0,99 & 1,25$\times$ & 1,36$\times$ & 40,1 & 40,2 & 1,63 & 1,69 & 0,00 \\
     & BN-Scale & 64,23 & +1,85 & 0,97 & 1,28$\times$ & 1,53$\times$ & 37,1 & 48,6 & 1,69 & 1,77 & 0,00 \\
     & FPGM & 64,64 & +2,26 & 1,00 & 1,24$\times$ & 1,21$\times$ & 38,9 & 29,3 & 1,65 & 1,69 & 0,00 \\
     & Taylor & 64,60 & +2,22 & 0,97 & 1,28$\times$ & 1,39$\times$ & 39,8 & 44,5 & 1,63 & 1,70 & 0,00 \\
     & Random & 61,79 & -0,59 & 1,15 & 1,08$\times$ & 1,25$\times$ & 40,1 & 41,7 & 1,71 & 1,83 & 0,00 \\
    \midrule
    \multirow{6}{*}{50\%} & Magnitude L1 & 64,69 & +2,31 & 0,91 & 1,36$\times$ & 1,51$\times$ & 50,0 & 49,8 & 1,36 & 1,41 & 0,00 \\
     & Magnitude L2 & 64,08 & +1,70 & 0,93 & 1,33$\times$ & 1,41$\times$ & 50,0 & 47,2 & 1,36 & 1,41 & 0,00 \\
     & BN-Scale & 64,64 & +2,26 & 0,91 & 1,37$\times$ & 1,65$\times$ & 44,3 & 56,6 & 1,49 & 1,55 & 0,00 \\
     & FPGM & 64,30 & +1,92 & 0,90 & 1,37$\times$ & 1,34$\times$ & 61,0 & 41,9 & 1,09 & 1,12 & 0,00 \\
     & Taylor & 64,96 & +2,58 & 0,91 & 1,36$\times$ & 1,55$\times$ & 50,0 & 51,6 & 1,36 & 1,41 & 0,00 \\
     & Random & 61,09 & -1,29 & 1,04 & 1,19$\times$ & 1,42$\times$ & 50,1 & 52,5 & 1,44 & 1,52 & 0,00 \\
    \midrule
    \multirow{5}{*}{60\%} & Magnitude L1 & 64,12 & +1,74 & 0,87 & 1,43$\times$ & 1,68$\times$ & 60,1 & 59,0 & 1,09 & 1,14 & 0,00 \\
     & Magnitude L2 & 64,42 & +2,04 & 0,87 & 1,43$\times$ & 1,69$\times$ & 60,0 & 59,1 & 1,09 & 1,14 & 0,00 \\
     & BN-Scale & 64,42 & +2,04 & 0,85 & 1,46$\times$ & 1,74$\times$ & 55,9 & 61,0 & 1,19 & 1,23 & 0,00 \\
     & Taylor & 64,57 & +2,19 & 0,87 & 1,43$\times$ & 1,70$\times$ & 60,0 & 59,6 & 1,09 & 1,12 & 0,00 \\
     & Random & 58,99 & -3,39 & 0,99 & 1,25$\times$ & 1,62$\times$ & 60,4 & 59,9 & 1,18 & 1,31 & 0,00 \\
    \midrule
    \multirow{6}{*}{70\%} & Magnitude L1 & 64,66 & +2,28 & 0,80 & 1,55$\times$ & 1,90$\times$ & 70,2 & 69,4 & 0,83 & 0,88 & 0,00 \\
     & Magnitude L2 & 63,94 & +1,56 & 0,80 & 1,56$\times$ & 1,90$\times$ & 70,7 & 69,7 & 0,82 & 0,86 & 0,00 \\
     & BN-Scale & 64,34 & +1,96 & 0,80 & 1,55$\times$ & 1,88$\times$ & 70,3 & 70,5 & 0,83 & 0,89 & 0,00 \\
     & FPGM & 64,85 & +2,47 & 0,80 & 1,56$\times$ & 1,85$\times$ & 70,1 & 67,2 & 0,84 & 0,87 & 0,00 \\
     & Taylor & 64,60 & +2,22 & 0,78 & 1,58$\times$ & 2,01$\times$ & 70,2 & 70,5 & 0,83 & 0,86 & 0,00 \\
     & Random & 56,34 & -6,04 & 0,92 & 1,35$\times$ & 1,96$\times$ & 70,2 & 70,5 & 0,92 & 1,02 & 0,00 \\
    \midrule
    \multirow{6}{*}{80\%} & Magnitude L1 & 64,83 & +2,45 & 0,77 & 1,61$\times$ & 2,11$\times$ & 80,2 & 74,0 & 0,60 & 0,65 & 0,00 \\
     & Magnitude L2 & 64,09 & +1,71 & 0,78 & 1,59$\times$ & 2,13$\times$ & 80,3 & 73,9 & 0,60 & 0,65 & 0,00 \\
     & BN-Scale & 63,98 & +1,60 & 0,77 & 1,61$\times$ & 2,24$\times$ & 80,3 & 76,2 & 0,60 & 0,64 & 0,00 \\
     & FPGM & 64,36 & +1,98 & 0,77 & 1,61$\times$ & 2,18$\times$ & 80,1 & 73,0 & 0,61 & 0,64 & 0,00 \\
     & Taylor & 64,13 & +1,75 & 0,76 & 1,62$\times$ & 2,23$\times$ & 80,2 & 74,6 & 0,59 & 0,64 & 0,00 \\
     & Random & -- & -- & -- & -- & 2,57$\times$ & 80,2 & 80,7 & 0,65 & 0,00 & 0,00 \\
    \midrule
    \multirow{6}{*}{90\%} & Magnitude L1 & 51,76 & -10,62 & 0,70 & 1,76$\times$ & 4,40$\times$ & 90,1 & 89,8 & 0,35 & 0,37 & 0,00 \\
     & Magnitude L2 & 55,82 & -6,56 & 0,71 & 1,74$\times$ & 4,59$\times$ & 90,1 & 90,7 & 0,35 & 0,37 & 0,00 \\
     & BN-Scale & 59,09 & -3,29 & 0,73 & 1,70$\times$ & 3,12$\times$ & 90,2 & 89,0 & 0,35 & 0,39 & 0,00 \\
     & FPGM & 55,75 & -6,63 & 0,69 & 1,81$\times$ & 4,39$\times$ & 90,1 & 91,3 & 0,34 & 0,37 & 0,00 \\
     & Taylor & 61,94 & -0,44 & 0,73 & 1,69$\times$ & 2,74$\times$ & 90,1 & 82,0 & 0,34 & 0,38 & 0,00 \\
     & Random & 40,06 & -22,32 & 0,75 & 1,65$\times$ & 3,49$\times$ & 90,2 & 90,5 & 0,38 & 0,43 & 0,00 \\
    \bottomrule
  \end{tabular}}}
\end{table}

\clearpage
\renewcommand{\benchscale}{0.70}
\begin{table}[H]
  \centering
  \caption{Mesures complètes de GoogLeNet élagué, en INT8.}
  \label{tab:bench-googlenet}
  \centerline{\resizebox{\benchscale\linewidth}{!}{%
  \begin{tabular}{c l rr rrr rr rrr}
    \toprule
    & & \multicolumn{2}{c}{Précision} & \multicolumn{3}{c}{Temps} & \multicolumn{2}{c}{Réduction (\%)} & \multicolumn{3}{c}{Mémoire (Mio)} \\
    \cmidrule(lr){3-4} \cmidrule(lr){5-7} \cmidrule(lr){8-9} \cmidrule(lr){10-12}
    Cible & Critère & Top-1 & $\Delta$ Top-1 & Lat. (ms) & Accél. TPU & Accél. CPU & Param. & MACs & Taille & Interne & Externe \\
    \midrule
    -- & référence & 77,02 & +0,00 & 2,04 & 1,00$\times$ & 1,00$\times$ & 0,0 & 0,0 & 5,66 & 5,94 & 0,00 \\
    \midrule
    \multirow{7}{*}{10\%} & Magnitude L1 & 77,56 & +0,54 & 2,21 & 0,92$\times$ & 1,12$\times$ & 10,0 & 17,3 & 5,10 & 5,44 & 0,00 \\
     & Magnitude L2 & 77,24 & +0,22 & 2,17 & 0,94$\times$ & 1,13$\times$ & 10,0 & 17,1 & 5,10 & 5,43 & 0,00 \\
     & BN-Scale & 77,44 & +0,42 & 1,95 & 1,04$\times$ & 1,18$\times$ & 10,1 & 21,2 & 5,10 & 5,55 & 0,00 \\
     & FPGM & 77,42 & +0,40 & 2,10 & 0,97$\times$ & 1,20$\times$ & 10,0 & 22,6 & 5,10 & 5,42 & 0,00 \\
     & Taylor & 77,40 & +0,38 & 1,99 & 1,02$\times$ & 1,11$\times$ & 10,0 & 15,2 & 5,10 & 5,54 & 0,00 \\
     & OBD-C & 77,45 & +0,43 & 1,96 & 1,04$\times$ & 1,09$\times$ & 10,0 & 13,4 & 5,10 & 5,52 & 0,00 \\
     & Random & 77,65 & +0,63 & 2,12 & 0,96$\times$ & 1,06$\times$ & 10,0 & 10,0 & 5,11 & 5,69 & 0,00 \\
    \midrule
    \multirow{7}{*}{20\%} & Magnitude L1 & 77,58 & +0,56 & 1,87 & 1,09$\times$ & 1,38$\times$ & 20,0 & 34,6 & 4,55 & 5,02 & 0,00 \\
     & Magnitude L2 & 76,82 & -0,20 & 1,87 & 1,09$\times$ & 1,46$\times$ & 20,1 & 40,0 & 4,55 & 5,01 & 0,00 \\
     & BN-Scale & 77,32 & +0,30 & 1,85 & 1,10$\times$ & 1,43$\times$ & 20,1 & 37,9 & 4,55 & 5,05 & 0,00 \\
     & FPGM & 77,12 & +0,10 & 1,85 & 1,10$\times$ & 1,40$\times$ & 20,1 & 36,6 & 4,55 & 5,00 & 0,00 \\
     & Taylor & 77,04 & +0,02 & 2,04 & 1,00$\times$ & 1,17$\times$ & 20,0 & 20,7 & 4,55 & 5,02 & 0,00 \\
     & OBD-C & 77,23 & +0,21 & 2,13 & 0,96$\times$ & 1,21$\times$ & 20,0 & 23,7 & 4,55 & 5,10 & 0,00 \\
     & Random & 77,54 & +0,52 & 2,23 & 0,92$\times$ & 1,16$\times$ & 20,1 & 19,8 & 4,55 & 5,34 & 0,00 \\
    \midrule
    \multirow{7}{*}{30\%} & Magnitude L1 & 76,78 & -0,24 & 1,60 & 1,27$\times$ & 1,80$\times$ & 30,1 & 53,9 & 3,99 & 4,50 & 0,00 \\
     & Magnitude L2 & 76,41 & -0,61 & 1,71 & 1,19$\times$ & 1,82$\times$ & 30,0 & 54,4 & 3,99 & 4,49 & 0,00 \\
     & BN-Scale & 76,77 & -0,25 & 1,67 & 1,22$\times$ & 1,71$\times$ & 30,1 & 50,2 & 3,99 & 4,49 & 0,00 \\
     & FPGM & 77,08 & +0,06 & 1,81 & 1,13$\times$ & 1,69$\times$ & 30,1 & 49,0 & 3,99 & 4,35 & 0,00 \\
     & Taylor & 77,19 & +0,17 & 2,06 & 0,99$\times$ & 1,26$\times$ & 30,1 & 26,9 & 3,99 & 4,66 & 0,00 \\
     & OBD-C & 77,23 & +0,21 & 2,15 & 0,95$\times$ & 1,29$\times$ & 30,2 & 28,4 & 3,99 & 4,71 & 0,00 \\
     & Random & 76,81 & -0,21 & 2,08 & 0,98$\times$ & 1,29$\times$ & 30,1 & 30,1 & 3,99 & 4,79 & 0,00 \\
    \midrule
    \multirow{7}{*}{40\%} & Magnitude L1 & 75,98 & -1,04 & 1,63 & 1,25$\times$ & 2,31$\times$ & 40,0 & 65,2 & 3,43 & 3,96 & 0,00 \\
     & Magnitude L2 & 75,69 & -1,33 & 1,47 & 1,39$\times$ & 2,35$\times$ & 40,1 & 65,8 & 3,43 & 3,95 & 0,00 \\
     & BN-Scale & 76,04 & -0,98 & 1,54 & 1,33$\times$ & 2,13$\times$ & 40,1 & 61,9 & 3,43 & 4,00 & 0,00 \\
     & FPGM & 76,70 & -0,32 & 1,71 & 1,19$\times$ & 2,12$\times$ & 40,1 & 62,2 & 3,43 & 3,90 & 0,00 \\
     & Taylor & 76,69 & -0,33 & 1,79 & 1,14$\times$ & 1,39$\times$ & 40,1 & 34,6 & 3,44 & 4,14 & 0,00 \\
     & OBD-C & 77,03 & +0,01 & 2,04 & 1,00$\times$ & 1,38$\times$ & 40,0 & 34,5 & 3,44 & 4,15 & 0,00 \\
     & Random & 76,45 & -0,57 & 1,83 & 1,12$\times$ & 1,52$\times$ & 40,2 & 40,7 & 3,44 & 4,15 & 0,00 \\
    \midrule
    \multirow{7}{*}{50\%} & Magnitude L1 & 75,30 & -1,72 & 1,47 & 1,38$\times$ & 2,90$\times$ & 50,0 & 73,6 & 2,87 & 3,43 & 0,00 \\
     & Magnitude L2 & 75,19 & -1,83 & 1,42 & 1,44$\times$ & 2,91$\times$ & 50,2 & 74,1 & 2,87 & 3,39 & 0,00 \\
     & BN-Scale & 75,43 & -1,59 & 1,59 & 1,29$\times$ & 2,71$\times$ & 50,2 & 71,5 & 2,87 & 3,51 & 0,00 \\
     & FPGM & 76,07 & -0,95 & 1,50 & 1,36$\times$ & 2,80$\times$ & 50,1 & 73,0 & 2,87 & 3,36 & 0,00 \\
     & Taylor & 76,09 & -0,93 & 1,99 & 1,02$\times$ & 1,50$\times$ & 50,1 & 40,6 & 2,89 & 3,58 & 0,00 \\
     & OBD-C & 76,30 & -0,72 & 1,62 & 1,26$\times$ & 1,52$\times$ & 50,1 & 41,0 & 2,88 & 3,62 & 0,00 \\
     & Random & 76,72 & -0,30 & 1,76 & 1,16$\times$ & 1,77$\times$ & 50,1 & 51,3 & 2,88 & 3,62 & 0,00 \\
    \midrule
    \multirow{7}{*}{60\%} & Magnitude L1 & 73,54 & -3,48 & 1,25 & 1,64$\times$ & 4,04$\times$ & 60,0 & 82,7 & 2,31 & 2,84 & 0,00 \\
     & Magnitude L2 & 73,80 & -3,22 & 1,51 & 1,35$\times$ & 3,98$\times$ & 60,2 & 82,5 & 2,30 & 2,81 & 0,00 \\
     & BN-Scale & 74,29 & -2,73 & 1,41 & 1,45$\times$ & 3,37$\times$ & 60,1 & 78,8 & 2,31 & 2,95 & 0,00 \\
     & FPGM & 73,11 & -3,91 & 1,23 & 1,66$\times$ & 3,70$\times$ & 60,1 & 80,7 & 2,31 & 2,70 & 0,00 \\
     & Taylor & 75,65 & -1,37 & 1,81 & 1,13$\times$ & 1,70$\times$ & 60,3 & 48,0 & 2,32 & 2,97 & 0,00 \\
     & OBD-C & 75,52 & -1,50 & 1,75 & 1,16$\times$ & 1,82$\times$ & 60,0 & 52,6 & 2,33 & 2,96 & 0,00 \\
     & Random & 76,12 & -0,90 & 1,52 & 1,34$\times$ & 2,00$\times$ & 60,2 & 58,6 & 2,32 & 2,98 & 0,00 \\
    \midrule
    \multirow{7}{*}{70\%} & Magnitude L1 & 72,36 & -4,66 & 1,10 & 1,85$\times$ & 5,03$\times$ & 70,2 & 87,0 & 1,75 & 2,26 & 0,00 \\
     & Magnitude L2 & 71,08 & -5,94 & 1,13 & 1,80$\times$ & 5,18$\times$ & 70,0 & 87,4 & 1,75 & 2,16 & 0,00 \\
     & BN-Scale & 71,76 & -5,26 & 1,26 & 1,62$\times$ & 4,50$\times$ & 70,0 & 85,2 & 1,76 & 2,22 & 0,00 \\
     & FPGM & 70,94 & -6,08 & 1,14 & 1,78$\times$ & 5,04$\times$ & 70,2 & 87,1 & 1,74 & 2,24 & 0,00 \\
     & Taylor & 74,57 & -2,45 & 1,62 & 1,26$\times$ & 1,91$\times$ & 70,1 & 53,9 & 1,77 & 2,37 & 0,00 \\
     & OBD-C & 74,75 & -2,27 & 1,46 & 1,40$\times$ & 2,10$\times$ & 70,2 & 60,1 & 1,76 & 2,33 & 0,00 \\
     & Random & 75,50 & -1,52 & 1,41 & 1,45$\times$ & 2,56$\times$ & 70,1 & 69,1 & 1,77 & 2,37 & 0,00 \\
    \midrule
    \multirow{7}{*}{80\%} & Magnitude L1 & 70,12 & -6,90 & 1,00 & 2,03$\times$ & 7,13$\times$ & 80,3 & 91,7 & 1,18 & 1,51 & 0,00 \\
     & Magnitude L2 & 69,02 & -8,00 & 1,03 & 1,98$\times$ & 7,16$\times$ & 80,0 & 91,8 & 1,19 & 1,55 & 0,00 \\
     & BN-Scale & 67,60 & -9,42 & 1,05 & 1,94$\times$ & 6,55$\times$ & 80,3 & 91,8 & 1,18 & 1,55 & 0,00 \\
     & FPGM & 64,56 & -12,46 & 1,05 & 1,94$\times$ & 7,42$\times$ & 80,0 & 92,1 & 1,19 & 1,56 & 0,00 \\
     & Taylor & 73,96 & -3,06 & 1,64 & 1,25$\times$ & 2,34$\times$ & 80,1 & 64,8 & 1,21 & 1,70 & 0,00 \\
     & OBD-C & 73,46 & -3,56 & 1,48 & 1,38$\times$ & 2,70$\times$ & 80,1 & 70,6 & 1,20 & 1,68 & 0,00 \\
     & Random & 75,01 & -2,01 & 1,32 & 1,54$\times$ & 3,39$\times$ & 80,1 & 78,9 & 1,20 & 1,67 & 0,00 \\
    \midrule
    \multirow{7}{*}{90\%} & Magnitude L1 & 65,89 & -11,13 & 0,94 & 2,17$\times$ & 10,89$\times$ & 90,1 & 95,7 & 0,62 & 0,89 & 0,00 \\
     & Magnitude L2 & 63,78 & -13,24 & 0,92 & 2,21$\times$ & 10,68$\times$ & 90,2 & 96,0 & 0,62 & 0,88 & 0,00 \\
     & BN-Scale & 63,15 & -13,87 & 0,93 & 2,19$\times$ & 10,19$\times$ & 90,0 & 95,9 & 0,63 & 0,89 & 0,00 \\
     & FPGM & 63,17 & -13,85 & 0,93 & 2,19$\times$ & 11,60$\times$ & 90,0 & 96,1 & 0,63 & 0,90 & 0,00 \\
     & Taylor & 71,10 & -5,92 & 1,26 & 1,62$\times$ & 3,06$\times$ & 90,0 & 75,0 & 0,64 & 1,04 & 0,00 \\
     & OBD-C & 70,42 & -6,60 & 1,17 & 1,74$\times$ & 4,24$\times$ & 90,0 & 83,6 & 0,64 & 1,00 & 0,00 \\
     & Random & 71,45 & -5,57 & 1,04 & 1,97$\times$ & 5,75$\times$ & 90,1 & 89,8 & 0,63 & 0,93 & 0,00 \\
    \bottomrule
  \end{tabular}}}
\end{table}

\clearpage
\renewcommand{\benchscale}{0.88}
\begin{table}[H]
  \centering
  \caption{Mesures complètes de SqueezeNet 1.1 élagué, en INT8.}
  \label{tab:bench-squeezenet1_1}
  \centerline{\resizebox{\benchscale\linewidth}{!}{%
  \begin{tabular}{c l rr rrr rr rrr}
    \toprule
    & & \multicolumn{2}{c}{Précision} & \multicolumn{3}{c}{Temps} & \multicolumn{2}{c}{Réduction (\%)} & \multicolumn{3}{c}{Mémoire (Mio)} \\
    \cmidrule(lr){3-4} \cmidrule(lr){5-7} \cmidrule(lr){8-9} \cmidrule(lr){10-12}
    Cible & Critère & Top-1 & $\Delta$ Top-1 & Lat. (ms) & Accél. TPU & Accél. CPU & Param. & MACs & Taille & Interne & Externe \\
    \midrule
    -- & référence & 50,21 & +0,00 & 0,67 & 1,00$\times$ & 1,00$\times$ & 0,0 & 0,0 & 0,83 & 0,82 & 0,00 \\
    \midrule
    \multirow{5}{*}{10\%} & Magnitude L1 & 54,68 & +4,47 & 0,68 & 0,98$\times$ & 1,01$\times$ & 10,0 & 9,2 & 0,76 & 0,78 & 0,00 \\
     & Magnitude L2 & 53,40 & +3,19 & 0,69 & 0,97$\times$ & 1,01$\times$ & 10,0 & 9,2 & 0,75 & 0,78 & 0,00 \\
     & FPGM & 53,58 & +3,37 & 0,66 & 1,01$\times$ & 1,10$\times$ & 6,9 & 6,6 & 0,77 & 0,79 & 0,00 \\
     & Taylor & 54,24 & +4,03 & 0,69 & 0,96$\times$ & 1,00$\times$ & 10,0 & 9,6 & 0,76 & 0,78 & 0,00 \\
     & Random & 53,13 & +2,92 & 0,68 & 0,98$\times$ & 1,00$\times$ & 10,1 & 10,7 & 0,75 & 0,79 & 0,00 \\
    \midrule
    \multirow{5}{*}{20\%} & Magnitude L1 & 54,24 & +4,03 & 0,65 & 1,03$\times$ & 1,25$\times$ & 20,0 & 21,1 & 0,67 & 0,73 & 0,00 \\
     & Magnitude L2 & 54,68 & +4,47 & 0,67 & 1,00$\times$ & 1,24$\times$ & 20,1 & 20,9 & 0,67 & 0,73 & 0,00 \\
     & FPGM & 54,13 & +3,92 & 0,64 & 1,05$\times$ & 1,23$\times$ & 15,9 & 14,7 & 0,70 & 0,71 & 0,00 \\
     & Taylor & 53,15 & +2,94 & 0,65 & 1,02$\times$ & 1,19$\times$ & 20,0 & 21,8 & 0,67 & 0,73 & 0,00 \\
     & Random & 53,99 & +3,78 & 0,69 & 0,96$\times$ & 1,06$\times$ & 20,2 & 15,5 & 0,68 & 0,74 & 0,00 \\
    \midrule
    \multirow{5}{*}{30\%} & Magnitude L1 & 52,14 & +1,93 & 0,63 & 1,05$\times$ & 1,46$\times$ & 30,1 & 30,6 & 0,59 & 0,65 & 0,00 \\
     & Magnitude L2 & 52,99 & +2,78 & 0,63 & 1,06$\times$ & 1,49$\times$ & 30,5 & 31,1 & 0,59 & 0,65 & 0,00 \\
     & FPGM & 53,39 & +3,18 & 0,63 & 1,06$\times$ & 1,35$\times$ & 28,6 & 25,6 & 0,60 & 0,65 & 0,00 \\
     & Taylor & 52,79 & +2,58 & 0,66 & 1,01$\times$ & 1,34$\times$ & 30,0 & 31,2 & 0,59 & 0,68 & 0,00 \\
     & Random & 53,67 & +3,46 & 0,65 & 1,03$\times$ & 1,16$\times$ & 30,0 & 26,8 & 0,60 & 0,71 & 0,00 \\
    \midrule
    \multirow{5}{*}{40\%} & Magnitude L1 & 52,04 & +1,83 & 0,62 & 1,07$\times$ & 1,58$\times$ & 40,1 & 39,5 & 0,51 & 0,61 & 0,00 \\
     & Magnitude L2 & 51,87 & +1,66 & 0,62 & 1,07$\times$ & 1,62$\times$ & 40,1 & 38,3 & 0,51 & 0,59 & 0,00 \\
     & FPGM & 53,38 & +3,17 & 0,64 & 1,04$\times$ & 1,57$\times$ & 40,1 & 38,4 & 0,51 & 0,60 & 0,00 \\
     & Taylor & 52,47 & +2,26 & 0,63 & 1,06$\times$ & 1,46$\times$ & 40,0 & 37,5 & 0,51 & 0,61 & 0,00 \\
     & Random & 50,57 & +0,36 & 0,65 & 1,03$\times$ & 1,33$\times$ & 40,1 & 42,9 & 0,52 & 0,65 & 0,00 \\
    \midrule
    \multirow{5}{*}{50\%} & Magnitude L1 & 50,61 & +0,40 & 0,61 & 1,09$\times$ & 1,87$\times$ & 50,2 & 49,8 & 0,43 & 0,50 & 0,00 \\
     & Magnitude L2 & 51,25 & +1,04 & 0,62 & 1,08$\times$ & 1,87$\times$ & 50,1 & 49,9 & 0,43 & 0,50 & 0,00 \\
     & FPGM & 51,83 & +1,62 & 0,61 & 1,09$\times$ & 1,93$\times$ & 50,1 & 53,5 & 0,43 & 0,50 & 0,00 \\
     & Taylor & 53,41 & +3,20 & 0,63 & 1,05$\times$ & 1,58$\times$ & 50,1 & 44,9 & 0,43 & 0,52 & 0,00 \\
     & Random & 52,01 & +1,80 & 0,65 & 1,03$\times$ & 1,42$\times$ & 50,2 & 46,3 & 0,44 & 0,53 & 0,00 \\
    \midrule
    \multirow{5}{*}{60\%} & Magnitude L1 & 48,83 & -1,38 & 0,61 & 1,09$\times$ & 2,45$\times$ & 60,1 & 64,8 & 0,35 & 0,44 & 0,00 \\
     & Magnitude L2 & 46,15 & -4,06 & 0,61 & 1,09$\times$ & 2,50$\times$ & 59,6 & 63,4 & 0,35 & 0,44 & 0,00 \\
     & FPGM & 49,72 & -0,49 & 0,60 & 1,12$\times$ & 2,47$\times$ & 60,0 & 65,1 & 0,35 & 0,43 & 0,00 \\
     & Taylor & 53,11 & +2,90 & 0,63 & 1,07$\times$ & 1,70$\times$ & 60,2 & 50,0 & 0,35 & 0,45 & 0,00 \\
     & Random & 46,53 & -3,68 & 0,62 & 1,07$\times$ & 1,63$\times$ & 60,4 & 58,4 & 0,36 & 0,44 & 0,00 \\
    \midrule
    \multirow{5}{*}{70\%} & Magnitude L1 & 46,35 & -3,86 & 0,59 & 1,14$\times$ & 2,90$\times$ & 70,4 & 73,4 & 0,27 & 0,33 & 0,00 \\
     & Magnitude L2 & 39,83 & -10,38 & 0,59 & 1,13$\times$ & 2,89$\times$ & 70,1 & 73,0 & 0,27 & 0,33 & 0,00 \\
     & FPGM & 42,03 & -8,18 & 0,58 & 1,14$\times$ & 3,27$\times$ & 70,2 & 79,5 & 0,27 & 0,33 & 0,00 \\
     & Taylor & 51,37 & +1,16 & 0,63 & 1,05$\times$ & 1,81$\times$ & 70,3 & 55,8 & 0,27 & 0,38 & 0,00 \\
     & Random & 36,48 & -13,73 & 0,60 & 1,11$\times$ & 1,93$\times$ & 70,0 & 70,2 & 0,28 & 0,38 & 0,00 \\
    \midrule
    \multirow{5}{*}{80\%} & Magnitude L1 & 35,22 & -14,99 & 0,59 & 1,12$\times$ & 3,77$\times$ & 80,2 & 83,0 & 0,19 & 0,26 & 0,00 \\
     & Magnitude L2 & 26,76 & -23,45 & 0,57 & 1,16$\times$ & 4,00$\times$ & 80,4 & 85,3 & 0,19 & 0,25 & 0,00 \\
     & FPGM & 29,66 & -20,55 & 0,57 & 1,17$\times$ & 4,11$\times$ & 80,0 & 87,6 & 0,19 & 0,25 & 0,00 \\
     & Taylor & 40,99 & -9,22 & 0,61 & 1,09$\times$ & 1,98$\times$ & 80,1 & 63,0 & 0,19 & 0,24 & 0,00 \\
     & Random & 1,00 & -49,21 & 0,59 & 1,12$\times$ & 2,66$\times$ & 80,0 & 79,7 & 0,20 & 0,27 & 0,00 \\
    \midrule
    \multirow{5}{*}{90\%} & Magnitude L1 & 1,00 & -49,21 & 0,58 & 1,14$\times$ & 5,43$\times$ & 90,1 & 93,8 & 0,11 & 0,15 & 0,00 \\
     & Magnitude L2 & 1,00 & -49,21 & 0,57 & 1,18$\times$ & 5,05$\times$ & 90,2 & 92,6 & 0,11 & 0,15 & 0,00 \\
     & FPGM & 1,00 & -49,21 & 0,57 & 1,17$\times$ & 5,58$\times$ & 90,1 & 93,9 & 0,11 & 0,16 & 0,00 \\
     & Taylor & 27,87 & -22,34 & 0,60 & 1,11$\times$ & 2,79$\times$ & 90,1 & 79,4 & 0,11 & 0,16 & 0,00 \\
     & Random & 1,00 & -49,21 & 0,57 & 1,17$\times$ & 3,74$\times$ & 90,0 & 90,4 & 0,12 & 0,16 & 0,00 \\
    \bottomrule
  \end{tabular}}}
\end{table}

